% Source: https://artofproblemsolving.com/community/c3h2861454_mind_boggling
% Stable identity: aops-topic-2861454-q01
\textbf{Step 1 (an identity).} Fix $m\ge1$ and let $\theta_k=\dfrac{k\pi}{2m+1}$ for $k=1,\dots,m$. By de Moivre,
\[\sin((2m+1)\theta)=\operatorname{Im}(\cos\theta+i\sin\theta)^{2m+1}=\sin^{2m+1}\theta\sum_{j=0}^{m}(-1)^j\binom{2m+1}{2j+1}\cot^{2m-2j}\theta.\]
At $\theta=\theta_k$ the left side is $0$ and $\sin\theta_k\ne0$. So each $x_k=\cot^2\theta_k$ is a root of
\[p(x)=\binom{2m+1}{1}x^m-\binom{2m+1}{3}x^{m-1}+\cdots.\]
The $x_k$ are distinct (cotangent is decreasing and positive on $\left(0,\dfrac\pi2\right)$), so they are all $m$ roots of $p$. By Vieta,
\[\sum_{k=1}^m\cot^2\theta_k=\dfrac{\binom{2m+1}{3}}{\binom{2m+1}{1}}=\dfrac{m(2m-1)}{3}.\]

\textbf{Step 2 (an inequality).} For $0<\theta<\dfrac\pi2$, comparing triangle $OAP$, sector $OAP$ and triangle $OAT$ in the unit circle gives $\sin\theta<\theta<\tan\theta$.
\begin{asy}
usepackage("amsmath");
size(3.5cm);
defaultpen(fontsize(10pt));
real t=40;pair O=(0,0),A=(1,0),P=dir(t),Q=(1,tan(radians(t)));
draw(arc(O,1,-5,70));draw(O--A--Q--cycle);draw(O--P);draw(P--(P.x,0),dashed);
draw(arc(O,0.25,0,t));label("$\theta$",0.25*dir(t/2),E);
label("$O$",O,SW);label("$A$",A,SE);label("$P$",P,NW);label("$T$",Q,NE);
\end{asy}

Hence $\cot^2\theta<\dfrac1{\theta^2}<\dfrac{1}{\sin^2\theta}=1+\cot^2\theta$.

\textbf{Step 3 (squeeze).} Summing over $\theta=\theta_k$, and using $\dfrac1{\theta_k^2}=\dfrac{(2m+1)^2}{\pi^2k^2}$,
\[\dfrac{m(2m-1)}{3}<\dfrac{(2m+1)^2}{\pi^2}\sum_{k=1}^m\dfrac1{k^2}<m+\dfrac{m(2m-1)}3=\dfrac{2m(m+1)}{3}.\]
Multiply by $\dfrac{\pi^2}{(2m+1)^2}$ and let $m\to\infty$. Both $\dfrac{\pi^2m(2m-1)}{3(2m+1)^2}$ and $\dfrac{2\pi^2m(m+1)}{3(2m+1)^2}$ tend to $\dfrac{\pi^2}{6}$, so $\displaystyle\sum_{k=1}^\infty\dfrac1{k^2}=\dfrac{\pi^2}6$. $\blacksquare$
